\documentclass[12pt]{article}
\usepackage[utf8]{inputenc}
\usepackage[T1]{fontenc}
\usepackage{hyperref}
\title{Yola Language Resources for K-12 Curriculum Development}
\author{Compiled by Claude Sonnet 5.5 (claude-sonnet-5-5) for Andr\'as Kornai}
\date{}
\begin{document}
\maketitle

\section{Introduction}
Yola, also called the Forth and Bargy dialect, was a Middle English-derived variety spoken in County Wexford, Ireland. It became functionally extinct in the 19th century. Some sources claim a last speaker died in 1998, but this is contested. Some people have since tried to revive it. The search found a historical glossary and general articles, but no modern teaching materials, and it did not find who runs any revival effort.

\section{Primary Educational Resources}
None found. No primers, courses or school programmes turned up.

\section{Linguistic and Reference Materials}
\begin{itemize}
\item Overview of the dialect: \url{https://en.wikipedia.org/wiki/Yola_dialect}
\item The 1867 glossary, \emph{A Glossary, with some pieces of verse, of the old dialect of the English colony in the Baronies of Forth and Bargy, County of Wexford, Ireland}, edited by William Barnes from Jacob Poole's earlier work. It is described in the overview articles, but I did not see a direct link to it.
\item Encyclopedia article: \url{https://encyclopedia.pub/entry/27969}
\item Wexford Hub page on the dialect: \url{https://wexfordhub.com/yola-dialect/}
\end{itemize}

\section{Religious and Community Resources}
Yola continued to be used liturgically in Wexford churches well into the 20th century, according to one article in the results. No texts were found.

\section{Digital and Audio Resources}
\begin{itemize}
\item Endangered Languages Project record: \url{https://www.endangeredlanguages.com/node/110358/printable/print}
\end{itemize}
No audio or video found.

\section{Curriculum Development Context}
A heritage curriculum would have to be built almost entirely from the 1867 glossary and the verse and phrases in it, since no spoken tradition survives. A school could teach it as local history and dialect study, in connection with Wexford's heritage. Modern revival groups were not identified in the search, so the community side is unknown.
\end{document}